\section{Introduction}
A learned particle simulator builds its computation around an interaction graph: edges determine which particles exchange messages before the next position update \citep{sanchez2020}. Geometric neighborhoods provide a useful local prior, but they do not explicitly allocate additional interactions according to the current prediction problem. An omitted neighbor may contain information needed to predict a collision; elsewhere, more messages may add little or perturb a prediction that was already accurate. The challenge is to construct graph changes whose cost is controlled and to determine where those changes improve the simulation.

We study constructive adaptive interaction graphs: preserve the simulator's native graph, form a local annulus of optional pairs, and append a specified number of them. Every budgeted policy therefore adds the same number of pairs on a common particle state. Our residual controller ranks candidates using cached endpoint scores, requiring one simulator pass per subsequent forecast after initialization. We also train the simulator to encounter graph changes. A paired mixed-graph arm receives random budgeted expansions during training, while its base-only counterpart keeps the native graph; architecture, objective and sampled inputs remain matched. This intervention asks whether the model can learn to use additional messages before asking whether residual scores select useful ones.

These questions require different comparisons. Holding the evaluation policy fixed measures the effect of graph-exposure training. Holding the model and observed history fixed tests where a policy places its extra pairs. Autonomous rollouts then reveal the feedback between graph selection, predicted motion and future neighborhoods. A residual score estimates prediction error, which can reflect missing interactions, unobserved variation or model bias. Only some of these errors can be reduced by an added edge. Our analysis makes this distinction explicit through intervention benefit and conditional allocation bounds. It also separates the direct cost of expansion from geometry changes: fewer edges along one predicted trajectory need not mean cheaper allocation on the same input.

Exploratory studies with three paired seeds expose this separation empirically. Goop and Sand compare both training arms from initialization at 100k updates; WaterDrop compares paired 10k continuations of existing 100k models. Fixed-random rollout error improves in every Goop and WaterDrop seed. On Goop, its mean falls from $0.301$ to $0.129$, while observed-test residual allocation still loses to random in every mixed-arm seed (Figure~\ref{fig:goop-exposure-placement}). WaterDrop reverses that observed-test ordering, but its autonomous interaction has mixed seed signs. Sand is a stronger counterexample: exposure narrows the observed risk deficit yet worsens cached-risk rollouts in every seed, and the native base graph beats random expansion in both arms. These results motivate a method and evaluation framework that treats training support, placement quality and autonomous feedback separately, retaining adverse outcomes rather than equating residual prediction with useful computation.
